\chapter*{Lời nói đầu}
\addcontentsline{toc}{chapter}{Lời nói đầu} 
\markboth{Lời nói đầu}{} 

\noindent
Xin chào, mình là \textbf{Đinh Công Thái}, sinh viên lớp IT1-K67, ngành Khoa học Máy tính, Trường Đại học Bách khoa Hà Nội. Là nghiên cứu sinh tại Viên Nghiên cứu và Ứng dụng Trí tuệ nhân tạo của \textbf{PGS.TS Nguyễn Phi Lê}, đồng thời nghiên cứu và làm việc dưới sự hướng dẫn của anh \textbf{Nguyễn Tuấn Dũng}

\vspace{1em}
\noindent
\textbf{Vậy cuốn giáo trình này ra đời để làm gì?}  
Mình viết nó cho những bạn đang mơ hồ, không định hướng, không biết AI bắt đầu từ đâu, nhưng có chút tò mò và muốn thử. Nội dung tập trung vào \textbf{Thị giác máy tính (Computer Vision - CV)}.  

\vspace{1em}
\noindent
\textbf{Cấu trúc cuốn sách}  

\begin{itemize}[leftmargin=*]
	\item \textbf{Phần 1: Lý thuyết CV} – kể câu chuyện từ những nền tảng cơ bản cho tới các mô hình hiện đại như Transformer. Đọc phần này để biết “tại sao lại như vậy”.  
	\item \textbf{Phần 2: Thực chiến CV} – biến lý thuyết thành code, dự án, ứng dụng. Đọc phần này để biết “làm thế nào”.  
\end{itemize}

Hai phần này bổ sung cho nhau.

\begin{flushright}
	\textit{Trân trọng,} \\[0.5em]
	Đinh Công Thái
\end{flushright}